\chapter{HASIL DAN PEMBAHASAN}

Bab ini menyajikan hasil dari penelitian yang telah dilakukan serta
analisis terhadap hasil tersebut. Berbeda dengan Bab 3 yang menjelaskan
metodologi penelitian, Bab ini berfokus pada penyajian hasil eksperimen
serta pembahasan terhadap temuan yang diperoleh.

Hasil penelitian dapat disajikan dalam berbagai bentuk seperti tabel,
gambar, grafik, diagram, maupun visualisasi lainnya yang mampu
merepresentasikan hasil penelitian secara jelas.

Setiap tabel atau gambar yang ditampilkan harus dijelaskan dalam teks
dan dirujuk menggunakan perintah \texttt{\textbackslash ref}. Tidak
diperbolehkan menampilkan tabel atau gambar tanpa penjelasan.

\section{DATA}
\subsection{Cara Menyematkan Gambar}

Pada Bab ini, hasil penelitian sering direpresentasikan dalam bentuk gambar
seperti diagram, grafik, arsitektur sistem, maupun visualisasi lainnya.
Untuk menyematkan gambar pada dokumen \LaTeX, langkah-langkah yang perlu
dilakukan adalah sebagai berikut.

\subsubsection{Menambahkan Gambar ke Folder Project}

Sebelum menyematkan gambar pada dokumen, terlebih dahulu unggah file gambar
ke dalam folder \texttt{images} yang terdapat pada project \LaTeX.

Disarankan untuk menggunakan nama file yang jelas, menggunakan huruf kecil,
serta menghindari penggunaan spasi. Sebagai contoh:

\begin{verbatim}
images/diagram_sistem.png
images/alur_penelitian.png
images/hasil_eksperimen.png
\end{verbatim}

\subsubsection{Menulis Skrip Penyematan Gambar}

Setelah gambar ditambahkan ke dalam folder \texttt{images}, gambar dapat
disematkan menggunakan lingkungan \texttt{figure}. Contoh skrip dapat
dilihat sebagai berikut:

\begin{verbatim}
\begin{figure}[H]
\centering
\includegraphics[width=0.8\textwidth]{images/nama_gambar.png}
\caption{Judul gambar ditulis di sini}
\label{fig:nama-gambar}
\end{figure}
\end{verbatim}

Penjelasan bagian-bagian pada skrip di atas:

\begin{itemize}
\item \texttt{figure} : lingkungan untuk menampilkan gambar.
\item \texttt{centering} : memposisikan gambar di tengah halaman.
\item \texttt{includegraphics} : memanggil file gambar dari folder project.
\item \texttt{caption} : judul atau keterangan gambar.
\item \texttt{label} : penanda yang digunakan untuk memanggil gambar.
\end{itemize}

\subsubsection{Menyebutkan Gambar dalam Teks}

Setiap gambar yang ditampilkan harus dijelaskan dalam teks. Untuk
menyebutkan gambar, gunakan perintah \texttt{\textbackslash ref}. Contoh:

\begin{verbatim}
Seperti ditunjukkan pada Gambar \ref{fig:nama-gambar}, sistem yang
dikembangkan mampu melakukan proses klasifikasi secara otomatis.
\end{verbatim}

Dengan cara ini, nomor gambar akan dihasilkan secara otomatis oleh
\LaTeX sehingga penulis tidak perlu menuliskan nomor gambar secara
manual.

\subsection{Cara Menyematkan Tabel}

Selain gambar, hasil penelitian juga sering disajikan dalam bentuk tabel.
Tabel biasanya digunakan untuk menampilkan data numerik, hasil eksperimen,
maupun perbandingan metode.

\subsubsection{Menulis Skrip Tabel}

Untuk menyematkan tabel pada dokumen \LaTeX, gunakan lingkungan
\texttt{table} dan \texttt{tabular}. Contoh skrip tabel adalah sebagai
berikut:

\begin{verbatim}
\begin{table}[H]
\centering
\begin{tabular}{lccc}
\toprule
Metode & Accuracy & Precision & Recall \\
\midrule
Metode A & 0.85 & 0.83 & 0.82 \\
Metode B & 0.90 & 0.88 & 0.87 \\
Metode Usulan & 0.94 & 0.92 & 0.91 \\
\bottomrule
\end{tabular}
\caption{Hasil perbandingan performa metode}
\label{tab:hasil-metode}
\end{table}
\end{verbatim}

Penjelasan bagian skrip tabel:

\begin{itemize}
\item \texttt{table} : lingkungan untuk menampilkan tabel.
\item \texttt{tabular} : struktur tabel yang menentukan jumlah kolom.
\item \texttt{toprule}, \texttt{midrule}, dan \texttt{bottomrule} :
digunakan untuk membuat garis tabel yang lebih rapi.
\item \texttt{caption} : judul tabel.
\item \texttt{label} : penanda untuk referensi tabel.
\end{itemize}

\subsubsection{Menyebutkan Tabel dalam Teks}

Setiap tabel harus dijelaskan dalam teks sebelum atau sesudah tabel
ditampilkan. Untuk memanggil tabel gunakan perintah
\texttt{\textbackslash ref}. Contoh:

\begin{verbatim}
Hasil evaluasi metode dapat dilihat pada Tabel \ref{tab:hasil-metode}.
\end{verbatim}

Dengan cara ini, nomor tabel akan dihasilkan secara otomatis oleh
\LaTeX.

\subsection{Catatan Penting}

Dalam penyusunan Bab ini, setiap gambar atau tabel yang ditampilkan
harus selalu dijelaskan dalam teks. Tidak diperbolehkan menampilkan
gambar atau tabel tanpa pembahasan.

Contoh yang benar:

"Seperti ditunjukkan pada Gambar \ref{fig:research-flow}, sistem yang dikembangkan terdiri dari tiga komponen utama."

Contoh yang salah:

Menampilkan gambar atau tabel tanpa menjelaskan isi atau maknanya.

%%%%%%%%%%%%%%%%%%%%%%%%%%%%%%%%%%%%%%%%%%%%
\section{Hasil Implementasi Sistem}
%%%%%%%%%%%%%%%%%%%%%%%%%%%%%%%%%%%%%%%%%%%%

Bagian ini menjelaskan hasil implementasi dari sistem atau metode yang
telah dirancang pada Bab 3. Penjelasan dapat mencakup tampilan sistem,
proses implementasi, serta fitur yang dihasilkan.

%%%%%%%%%%%%%%%%%%%%%%%%%%%%%%%%%%%%%%%%%%%%
\section{Hasil Eksperimen}
%%%%%%%%%%%%%%%%%%%%%%%%%%%%%%%%%%%%%%%%%%%%

Bagian ini menyajikan hasil eksperimen yang dilakukan pada penelitian.
Hasil eksperimen biasanya disajikan dalam bentuk tabel atau grafik agar
lebih mudah dipahami.

\subsection{Hasil Pengujian Dataset}

Bagian ini menjelaskan hasil pengujian metode terhadap dataset yang
digunakan dalam penelitian.

%%%%%%%%%%%%%%%%%%%%%%%%%%%%%%%%%%%%%%%%%%%%
\section{Analisis Hasil Penelitian}
%%%%%%%%%%%%%%%%%%%%%%%%%%%%%%%%%%%%%%%%%%%%

Bagian ini berisi analisis terhadap hasil eksperimen yang telah
diperoleh. Penulis harus menjelaskan makna dari setiap hasil yang
ditampilkan serta menghubungkannya dengan tujuan penelitian.

\subsection{Perbandingan Metode}

Bagian ini menjelaskan perbandingan performa metode yang diusulkan
dengan metode lain yang telah ada sebelumnya.

%%%%%%%%%%%%%%%%%%%%%%%%%%%%%%%%%%%%%%%%%%%%
\section{Diskusi}
%%%%%%%%%%%%%%%%%%%%%%%%%%%%%%%%%%%%%%%%%%%%

Bagian diskusi berisi interpretasi hasil penelitian secara lebih
mendalam. Penulis dapat menjelaskan kelebihan metode yang diusulkan,
keterbatasan penelitian, serta implikasi hasil penelitian terhadap
bidang yang diteliti.